\section{E2E posts e comentários}
\label{sec:e2e-posts-comentarios}

Segunda vertical E2E permanente, derivada da Seção~9 (fluxos M12.1), da
Seção~8 (UI-11) e da Seção~7 (M12.1), executada contra o Firebase Emulator
Suite real e o frontend Next.js real. Reutiliza integralmente a infraestrutura
da Seção~\ref{sec:e2e-convites-auth} e o seed canônico
\texttt{functions/scripts/seed.ts}; nenhum segundo universo de fixtures foi
criado.

\subsection{Fixtures}

A turma \texttt{seed-turma-ultima-vaga} (Professor Alpha dono; Aluno
Matriculado como único membro) sustenta os cenários positivos e multiusuário. O
Professor Beta (dono de outra turma) e o Chefe Geral são usados nos negativos
de ownership. O aluno sem vínculo (\texttt{aluno.normal}) é usado no negativo
de alcance. A turma \texttt{seed-turma-colegas} (``Química Geral — T3
Colegas'', Professor Alpha dono) tem três membros --- o autor
\texttt{aluno.matriculado} e os colegas \texttt{aluno.colega.a} e
\texttt{aluno.colega.b} --- e sustenta o cenário de projeção COLEGA. Cada teste
cria seu próprio Post; não há dependência entre testes.

\subsection{Cenários}

\begin{longtable}{@{}p{0.14\textwidth}p{0.46\textwidth}p{0.32\textwidth}@{}}
\toprule
ID & Cenário & Camadas \\
\midrule
\endhead
POST-E2E-001 & Professor dono cria Post sem roteiro; o payload real do SDK
  omite \texttt{idRoteiroExperimento} e o documento persistido corresponde ao
  contrato & UI, callable \texttt{criarPost}, Firestore Emulator \\
POST-E2E-002 & Aluno participante vê o Post em contexto independente &
  2 identidades, Firestore Rules, \texttt{onSnapshot} \\
POST-E2E-003 & Chefe Geral (sem ownership) alcança o formulário, tenta criar e
  o backend rejeita; nada é persistido & UI, callable/autoridade, Firestore \\
POST-E2E-004 & Professor terceiro e aluno sem vínculo não alcançam a turma;
  nada é persistido & UI (painel lateral), Rules de list/get \\
COMMENT-E2E-001 & Aluno participante comenta e o comentário persiste com
  \texttt{id\_post}/\texttt{id\_usuario}/\texttt{texto} & UI, callable
  \texttt{adicionarComentario}, Firestore Emulator \\
COMMENT-E2E-002 & Professor dono observa o comentário do aluno &
  2 identidades, Firestore Rules \\
COMMENT-E2E-003 & Chefe Geral sem vínculo tenta comentar e o backend rejeita;
  nada é persistido & UI, callable/autoridade, Firestore \\
POST-E2E-005 & Professor dono remove Post com motivo; o aluno deixa de vê-lo e o
  documento permanece com \texttt{removido\_da\_apresentacao=true} &
  UI, callable \texttt{removerPost}, Firestore Emulator \\
COMMENT-E2E-004 & Professor dono modera comentário; o autor vê o original
  marcado e o auditor vê original + motivo & 2 identidades, callable
  \texttt{listarComentariosPost}, UI \\
ARQ-POST-E2E-001 & Turma arquivada não oferece superfície de escrita acadêmica;
  só ``Desarquivar'' no modal de arquivadas & UI (painel/modal), Firestore \\
COMMENT-E2E-005 & Colega de turma (turma de 3 alunos) recebe apenas aviso
  institucional e nunca o texto original do comentário moderado &
  3 identidades, callable \texttt{listarComentariosPost}, UI \\
\bottomrule
\end{longtable}

\subsection{Bug de implementação corrigido}

\textbf{Sintoma.} Publicar um Post \emph{sem} roteiro falhava com
\texttt{Erro de validação: idRoteiroExperimento: expected string, received
null}.

\textbf{Fonte normativa.} Seção~8 UI-11 e Seção~9, ``Publicar Post com e sem
roteiro'': o roteiro é opcional.

\textbf{Causa.} O serializador do SDK (\texttt{encode}) converte
\texttt{undefined} em \texttt{null}; \texttt{FeedTurma.tsx} enviava
\texttt{idRoteiroExperimento: idRoteiro || undefined}, e o schema Zod
(\texttt{z.string().optional()}) rejeita \texttt{null}. É a mesma classe do
defeito de \texttt{justificativaExcecao: null} da rodada anterior.

\textbf{Correção.} O produtor passa a omitir a chave quando não há roteiro, na
convenção já adotada em \texttt{convitesPayload.mjs}. Regressão permanente em
POST-E2E-001 (o payload real não contém a propriedade).

\subsection{Divergências conhecidas não mascaradas}

\begin{itemize}
  \item Edição de Post (M12.1) não possui callable nem UI (IMP-POST-001): o
    E2E de edição fica adiado até a implementação existir.
  \item Remoção lógica, moderação de comentários e leitura mascarada de
    Comentários foram implementadas nesta rodada (callables \texttt{removerPost},
    \texttt{moderarComentario}, \texttt{listarComentariosPost}), mas ainda carecem
    de E2E próprio; os negativos POST-E2E-003/COMMENT-E2E-003 agora provam a
    rejeição normativa de criação por Chefe\_Geral, que antes era um falso
    positivo aceito.
  \item O nome de exibição do autor cai no fallback ``Usuário''/``Professor''
    porque o seed não popula \texttt{nome} nas coleções de papel; a identidade
    correta é provada pela persistência (\texttt{id\_usuario}/
    \texttt{id\_professor}).
\end{itemize}

\subsection{Convergências desta rodada}

\begin{itemize}
  \item \textbf{Autorização M12.1:} criação de Post somente pelo professor-dono;
    criação de Comentário por professor-dono ou membro canônico com vínculo;
    Chefe\_Geral sem vínculo é rejeitado em ambos (Q13 não cria conteúdo).
  \item \textbf{M7/M9:} todos os callables de Posts/Comentários exigem
    \texttt{idOperacao}, revalidam autoridade persistida no commit e gravam
    receipt canônico.
  \item \textbf{Remoção lógica:} Post removido preserva documento, anexo e
    histórico, ocultando-se de colegas; motivo e autoria registrados.
  \item \textbf{Moderação:} comentário moderado preserva original, com histórico
    e projeção mascarada via endpoint autorizado; as três projeções são provadas
    em E2E (AUTOR/AUDITOR em COMMENT-E2E-004; COLEGA em COMMENT-E2E-005).
  \item \textbf{Security Rules:} leitura de Posts exige escopo acadêmico; Post
    removido só visível a dono/Chefe; Comentários tem leitura direta negada
    (acesso via \texttt{listarComentariosPost}).
  \item \textbf{Turma arquivada:} toda escrita acadêmica de Posts/Comentários é
    rejeitada sem efeito parcial.
\end{itemize}

\subsection{Testes de integração backend}

A vertical possui ainda evidência de integração backend em
\texttt{functions/src/\allowbreak integration/posts.integration.test.ts} (34
testes), exercitando os callables reais via \texttt{firebase-functions-test}
contra o Firestore Emulator: replay/reuso M7 (compatível não duplica;
incompatível é rejeitado), revogação/versão obsoleta M9, turma arquivada
rejeitando \texttt{criarPost}/\texttt{adicionarComentario}/\texttt{removerPost}/
\texttt{moderarComentario}, e as três projeções de leitura de Comentário
(AUTOR/COLEGA/AUDITOR), incluindo \texttt{texto=null} + aviso institucional
para o colega em comentário moderado. É nessa camada que a projeção COLEGA é
provada, complementando o E2E, que cobre AUTOR e AUDITOR.

\subsection{Camadas efetivamente exercitadas}

React/UI e Playwright, callables reais de Cloud Functions, Firestore Emulator e
Security Rules, Firebase Auth Emulator e múltiplas identidades em contextos
independentes. Nenhum mock de Firebase foi usado.
